\section{Introduction}
\label{sec:intro}

This document provides a template for academic papers based on the CVPR style format.
This template has been adapted to work as a general-purpose academic paper template, removing conference-specific submission requirements.

%-------------------------------------------------------------------------
\subsection{Language}

All manuscripts should be in English.

\subsection{Paper structure}
This template provides a standard academic paper structure with sections for introduction, methodology, experiments, results, and conclusion.
You can modify the section structure in the main.tex file by editing the input statements.

\subsection{Paper length}
For most academic venues, papers should follow the specific length requirements of your target conference or journal.
This template provides good formatting that can be adapted to various venues.

%-------------------------------------------------------------------------

\subsection{Mathematics}

Please number all of your sections and displayed equations as in these examples:
\begin{equation}
  E = m\cdot c^2
  \label{eq:important}
\end{equation}
and
\begin{equation}
  v = a\cdot t.
  \label{eq:also-important}
\end{equation}
It is important for the reader to be able to refer to any particular equation.
Just because you did not refer to it in the text does not mean that some future reader might not need to refer to it.
All authors will benefit from reading Mermin's description of how to write mathematics:
\url{http://www.pamitc.org/documents/mermin.pdf}.

\subsection{Citations and references}

When citing work, use proper academic citation format. 
For example, you can cite multiple works \cite{Alpher02,Alpher03,Alpher04}.
When citing a multi-author paper, you may save space by using ``et alia'', shortened to ``\etal'' (not ``{\em et.\ al.}'' as ``{\em et}'' is a complete word).
If you use the \verb'\etal' macro provided, then you need not worry about double periods when used at the end of a sentence as in Alpher \etal.
However, use it only when there are three or more authors.

\begin{figure}[t]
  \centering
  \fbox{\rule{0pt}{2in} \rule{0.9\linewidth}{0pt}}
   %\includegraphics[width=0.8\linewidth]{egfigure.eps}

   \caption{Example of caption.
   It is set in Roman so that mathematics (always set in Roman: $B \sin A = A \sin B$) may be included without an ugly clash.}
   \label{fig:onecol}
\end{figure}

\subsection{Figures and tables}

Include figures and tables as needed to support your research.
Use proper captions and labels so they can be referenced in the text.

\noindent
Compare the following:\\
\begin{tabular}{ll}
 \verb'$conf_a$' &  $conf_a$ \\
 \verb'$\mathit{conf}_a$' & $\mathit{conf}_a$
\end{tabular}\\
See The \TeX book, p165.

The space after \eg, meaning ``for example'', should not be a sentence-ending space.
So \eg is correct, {\em e.g.} is not.
The provided \verb'\eg' macro takes care of this.

\begin{figure*}
  \centering
  \begin{subfigure}{0.68\linewidth}
    \fbox{\rule{0pt}{2in} \rule{.9\linewidth}{0pt}}
    \caption{An example of a subfigure.}
    \label{fig:short-a}
  \end{subfigure}
  \hfill
  \begin{subfigure}{0.28\linewidth}
    \fbox{\rule{0pt}{2in} \rule{.9\linewidth}{0pt}}
    \caption{Another example of a subfigure.}
    \label{fig:short-b}
  \end{subfigure}
  \caption{Example of a short caption, which should be centered.}
  \label{fig:short}
\end{figure*}
